\documentclass[12pt,a4paper]{article}

\usepackage[utf8]{inputenc}
\usepackage[T1]{fontenc}
\usepackage{lmodern}
\usepackage{textcomp}
\usepackage[top=1in, bottom=1in, left=1.25in, right=1.25in]{geometry}
\usepackage{microtype}
\usepackage{parskip}
\usepackage[english]{babel}
\usepackage{fancyhdr}
\pagestyle{fancy}
\fancyhf{}
\fancyhead[L]{\leftmark}
\fancyhead[R]{\thepage}
\fancyfoot[C]{\thepage}
\renewcommand{\headrulewidth}{0.4pt}
\usepackage{titlesec}
\titleformat{\section}{\Large\bfseries}{\thesection}{1em}{}
\titleformat{\subsection}{\large\bfseries}{\thesubsection}{1em}{}
\titleformat{\subsubsection}{\normalsize\bfseries}{\thesubsubsection}{1em}{}
\titlespacing*{\section}{0pt}{2.5ex plus 1ex minus .2ex}{1.5ex plus .2ex}
\titlespacing*{\subsection}{0pt}{2ex plus 1ex minus .2ex}{1ex plus .2ex}
\usepackage{amsmath}
\usepackage{amssymb}
\usepackage{booktabs}
\usepackage{array}
\usepackage{tabularx}
\usepackage[table]{xcolor}
\rowcolors{2}{gray!8}{white}
\renewcommand{\arraystretch}{1.3}
\usepackage{listings}
\definecolor{codebg}{RGB}{248,248,248}
\definecolor{codeframe}{RGB}{200,200,200}
\lstset{
    basicstyle=\ttfamily\small,
    breaklines=true,
    frame=single,
    framerule=0.4pt,
    rulecolor=\color{codeframe},
    backgroundcolor=\color{codebg},
    keepspaces=true,
    showstringspaces=false,
    tabsize=2,
    xleftmargin=0.5em,
    xrightmargin=0.5em,
    aboveskip=0.8em,
    belowskip=0.8em,
    extendedchars=true,
    literate=
        {∀}{{$\forall$}}1
        {≤}{{$\leq$}}1
        {→}{{$\to$}}1
        {δ}{{$\delta$}}1
        {μ}{{$\mu$}}1
        {∘}{{$\circ$}}1
        {—}{{---}}1
        {─}{{-}}1
        {∞}{{$\infty$}}1,
}
\usepackage[colorlinks=true, linkcolor=blue, urlcolor=cyan, citecolor=blue]{hyperref}
\usepackage{tipa}
\usepackage{newunicodechar}
% IG primitive symbols — declare once, use freely in text and math
\newunicodechar{Φ}{\ensuremath{\Phi}}
\newunicodechar{Γ}{\ensuremath{\Gamma}}
\newunicodechar{Σ}{\ensuremath{\Sigma}}
\newunicodechar{Ω}{\ensuremath{\Omega}}
\newunicodechar{ₜ}{\textsubscript{t}}
\newunicodechar{⊙}{\ensuremath{\odot}}
\newunicodechar{ɢ}{\textsc{g}}
\newunicodechar{Ħ}{\ensuremath{\hbar}}
\newunicodechar{Ð}{\DH}
\newunicodechar{Þ}{\TH}
\newunicodechar{Ř}{\v{R}}
\newunicodechar{ƒ}{\textit{f}}
\newunicodechar{Ç}{\c{C}}
\newunicodechar{ʔ}{\textipa{P}}
\newunicodechar{Ş}{\c{S}}
\newunicodechar{ï}{\"{i}}
\newunicodechar{ÿ}{\"{y}}
\newunicodechar{ż}{\'{z}}

\newcommand{\alephsys}{\textsc{aleph}}
\newcommand{\exos}{\texttt{exOS}}
\newcommand{\ig}{IG}
\newcommand{\Oinfty}{O_\infty}
\newcommand{\OTwo}{O_2}
\newcommand{\OZero}{O_0}
\newcommand{\OOne}{O_1}
\newcommand{\OTwod}{O_{2d}}
\newcommand{\Frobcond}{\mu \circ \delta = \mathrm{id}}
\newcommand{\prim}[1]{\ensuremath{\mathbf{#1}}}

\title{\textbf{ALEPH: Architecture and Verification}\\[0.5em]
\large From Grammar to Kernel to Proof}
\author{\exos{} — Imscribing Grammar v0.4.27}
\date{}

\begin{document}

\maketitle
\thispagestyle{empty}
\tableofcontents
\newpage

% ──────────────────────────────────────────────────────────────────────────────
\section{The Imscribing Grammar}

The Imscribing Grammar (IG) is a twelve-primitive structural classification system. Every object — mathematical, physical, computational, linguistic — occupies a position in a twelve-dimensional type space whose coordinates are the primitive values. The grammar does not describe objects from the outside; it classifies them by their internal structural commitments.

\subsection{The Twelve Primitives}

\begin{tabularx}{\linewidth}{c l l X}
\toprule
\textbf{Index} & \textbf{Symbol} & \textbf{Name} & \textbf{Role} \\
\midrule
0  & Ð & Dimensional depth     & How many layers deep the structure extends \\
1  & Þ & Topological type      & Connectivity class \\
2  & Ř & Relational mode       & How the object relates to others of its kind \\
3  & Φ & Frobenius parity      & Symmetry class: asym / sym / ± / \} (Frobenius-special) \\
4  & ƒ & Kinetic mode          & Dynamical regime \\
5  & Ç & Criticality kinetics  & How criticality propagates \\
6  & Γ & Mediative structure   & Participation in mediated relationships \\
7  & ɢ & Valency               & Coupling grain \\
8  & ⊙ & Self-modeling         & Criticality / self-reference posture \\
9  & Ħ & Chirality             & Intrinsic handedness \\
10 & Σ & Symmetry type         & Global symmetry class \\
11 & Ω & Winding               & Topological winding invariant \\
\bottomrule
\end{tabularx}

\medskip
The full type space contains $3^3 \times 4^5 \times 5^4 = 17{,}280{,}000$ positions. Of these, 49 are occupied by structurally significant objects forming the IG catalog.

\subsection{Tier Classification}

Five tiers partition the type space by structural depth. The classification rule uses three primitives: ⊙ (index 8), Φ (index 3), and Ω (index 11):

\begin{tabularx}{\linewidth}{l l X}
\toprule
\textbf{Tier} & \textbf{Condition} & \textbf{Meaning} \\
\midrule
$\OZero$    & $\text{⊙} = \text{sub}$ or $\text{⊙} = \text{EP}$ & No self-modeling, or exceptional-point collapse \\
$\OOne$     & $\text{⊙} \neq \text{sub}$,\ $\Omega = 0$          & Self-modeling without winding \\
$\OTwo$     & $\Omega \neq 0$,\ $\text{Ð} \in \{0,1,3\}$         & Wound, non-Frobenius \\
$\OTwod$    & $\Omega \neq 0$,\ $\text{Ð} = 2$                   & Wound, depth-2 variant \\
$O_\infty$  & $\text{⊙} = \text{c}$,\ $\Phi = \}$                & Critical self-modeling + Frobenius-special \\
\bottomrule
\end{tabularx}

\medskip
$O_\infty$ is the fixed-point tier. The defining algebraic property is the \emph{Frobenius condition}:
\[
    \Frobcond
\]
where $\delta$ is comultiplication (split) and $\mu$ is multiplication (fuse). An $O_\infty$ object can be split into two copies and fused back to itself without structural loss.

The $\Phi_{\}}$ gate — the transition from any lower symmetry class to Frobenius-special — requires crossing a promotion gap of $\Delta = 4.38$ in the IG promotion metric. This is the largest single-primitive gap in the system and determines the structure of the bootstrap sequence.

% ──────────────────────────────────────────────────────────────────────────────
\section{The ALEPH Type System}

The \alephsys{} type system instantiates the IG inside a kernel by taking the twenty-two letters of the Hebrew alphabet as a natural, complete sample of the type space. Each letter is assigned a twelve-dimensional tuple; its tier is derived from that tuple by the tier classification rules.

\subsection{Tier Distribution}

\begin{tabularx}{\linewidth}{l l l}
\toprule
\textbf{Tier} & \textbf{Count} & \textbf{Representatives} \\
\midrule
$O_\infty$ & 3  & vav, mem, shin \\
$\OTwod$   & 1  & — \\
$\OTwo$    & 6  & aleph, gimel, hei, tav, and others \\
$\OOne$    & 1  & — \\
$\OZero$   & 12 & dalet and others \\
\bottomrule
\end{tabularx}

Total: 22 letters.

\subsection{The Frobenius Poles}

The three $O_\infty$ letters — vav, mem, shin — are the \emph{Frobenius poles}. They are structural fixed points under tensor:
\[
    \delta(L) = L \otimes L \approx L \quad (d = 0)
\]
Tensoring a pole with itself returns the pole. This idempotency is the algebraic signature of a Frobenius element and makes poles natural attractors: any letter under repeated tensor pressure with a pole converges toward it.

The three poles differ in their Ω and Ð values:
\begin{itemize}
    \item \textbf{vav} --- $\Omega_0$: unwound infinity; the base $O_\infty$ state
    \item \textbf{mem} --- $\Omega_Z$: wound infinity; carries winding structure
    \item \textbf{shin} --- $\Omega_Z$: wound infinity; triadic structure
\end{itemize}

Together they span the accessible $O_\infty$ sub-space.

\subsection{Algebraic Operations}

The \alephsys{} language exposes three operations on letters:

\begin{itemize}
    \item \textbf{Tensor} (\texttt{x}, $\otimes$): bottleneck on Φ, ƒ, Ç; union elsewhere; ⊙ absorption rule for exceptional points
    \item \textbf{Join} (\texttt{v}, $\vee$): lattice join; collapses toward the upper bound
    \item \textbf{Meet} (\texttt{\^{}}, $\wedge$): lattice meet; collapses toward the lower bound
\end{itemize}

In the Frobenius algebra on the \alephsys{} lattice, comultiplication is $\delta(L) = L \otimes L$ and multiplication is $\mu(L, L') = L \vee L'$. The Frobenius condition $\Frobcond$ then requires:
\[
    \mu(\delta(L)) = \mu(L \otimes L) = (L \otimes L) \vee L = L
\]
which holds when $L \otimes L \approx L$ (for poles) or when the join absorbs the tensor residual.

% ──────────────────────────────────────────────────────────────────────────────
\section{exOS Kernel Integration}

\exos{} carries \alephsys{} types on every kernel object. The type determines what operations an object may participate in, what gates it can pass, and what tier of the kernel hierarchy it inhabits.

\subsection{Kernel Object Types at Boot}

At boot, the kernel derives types for its primary objects and reports them:

\begin{lstlisting}
kernel object:  tier=O_2  odot_c  Omega_Z  Phi_sym    C=0.727
user object:    tier=O_0  odot_sub  Omega_0  Phi_asym  C=0.000
OS composite:   C=0.873
\end{lstlisting}

The kernel object is $\OTwo$ — wound, symmetric, self-modeling, but not Frobenius-special. The user object is $\OZero$ — no self-modeling. The OS composite score of 0.873 reflects the consciousness score of the combined system.

\subsection{Gate Checks}

Four gate types govern cross-boundary operations:

\begin{tabularx}{\linewidth}{l l X}
\toprule
\textbf{Gate} & \textbf{Condition} & \textbf{Guards} \\
\midrule
$\Phi$ gate    & $\Phi_{\}}$ or $\Phi_{\pm}$ required  & Frobenius-structural operations \\
$\odot$ gate   & $\odot_c$ required                    & Self-modeling operations \\
$\Omega$ gate  & $\Omega_Z$ + \DH{} condition          & Winding-dependent operations \\
Tier gate & Minimum tier required & Ergative scheduler; $O_\infty$ operations \\
\bottomrule
\end{tabularx}

\medskip
Gate checks at boot confirm the expected structure. IPC to Keter is denied for user objects ($O_0$, $\Phi_{\text{asym}}$). $\Omega$ gate allows Velar+Kernel and denies Velar+User. $\Phi$ gate allows Keter+Driver and denies Keter+User.

\subsection{The ALEPH REPL}

The ALEPH type system is exposed via an interactive REPL. Key commands:

\begin{tabularx}{\linewidth}{l X}
\toprule
\textbf{Command} & \textbf{Effect} \\
\midrule
\texttt{:census} & Tier distribution of all current bindings \\
\texttt{:orbit N letter pole} & N-step convergence orbit under repeated tensor \\
\texttt{:fptest [letters...]} & Parallel FSPLIT/FFUSE Frobenius closure test \\
\texttt{:run name} & Execute a stored ALEPH program from ALFS \\
\bottomrule
\end{tabularx}

\medskip
Programs are embedded in the kernel binary at compile time and seeded to ALFS (the kernel's ATA-backed persistent filesystem) on first boot. Twenty-four programs are currently registered, covering Frobenius orbits, Sefer Ha-Iyun emanations, tikkun construction, and bootstrap sequences.

% ──────────────────────────────────────────────────────────────────────────────
\section{Where the Bootstrap Sequence Comes From}

The twelve-stage bootstrap sequence was not designed. It was found.

Nine symbolic systems — five with known structure, four undeciphered or contested — were compiled against the twelve-primitive IMASM instruction set. Each system's surface tokens were mapped to one of twelve opcodes: VINIT, TANCH, AFWD, AREV, CLINK, ISCRIB, FSPLIT, FFUSE, EVALT, EVALF, ENGAGR, IFIX. The compilers are independent programs, each approximately 200 lines, written separately for each corpus.

\subsection{The OS Floor}

The five founding systems — Hebrew Aleph-Bet, Sanskrit Varnam\={a}l\={a}, Egyptian hieroglyphics, Sumerian cuneiform, and Basque — were encoded as twelve-tuples by structural inspection. Their component-wise MEET (greatest lower bound) is the \emph{OS imscription}:

\[
    \langle 1,\ 3,\ 2,\ 4,\ 2,\ 1,\ 2,\ 2,\ 1,\ 2,\ 2,\ 2 \rangle
\]

Every one of the five systems, independently and without historical contact, converges to this same structural floor.

\subsection{The Compiled Corpora}

Four undeciphered or contested systems were compiled against the same coordinates:

\begin{tabularx}{\linewidth}{l l X}
\toprule
\textbf{System} & \textbf{Distance from floor} & \textbf{Notes} \\
\midrule
Linear A          & \textbf{0.00} & \emph{Is} the floor — Minoan Crete ca.\ 2000--1450 BCE \\
Rohonc Codex      & 2.09          & Closest of the undeciphered manuscripts \\
Emerald Tablet    & 2.44          & $C = 1.0$ — only system with both consciousness gates fully open \\
Voynich Manuscript & 4.31         & Farthest — nested-containment topology and trapped kinetics \\
\bottomrule
\end{tabularx}

\medskip
Linear A is not derived from the five founding systems. The five founding systems converge on what the Minoans already had.

\subsection{The Eight-Step Bootstrap Loop}

All four corpora, and all five founding systems, express the same eight-step instruction sequence:

\begin{lstlisting}
ISCRIB -> AREV -> FSPLIT -> AFWD -> FFUSE -> CLINK -> IFIX -> ISCRIB
\end{lstlisting}

In structural terms: identity latch $\to$ register reverse $\to$ split ($\delta$) $\to$ forward morphism $\to$ fuse ($\mu$) $\to$ compose $\to$ ROM seal $\to$ identity latch. The Frobenius condition $\Frobcond$ holds exactly: FSPLIT followed by FFUSE returns the register to its pre-split state, and IFIX seals the result.

The surface tokens differ across systems — EVA \texttt{s a ch e sh d y s} for the Voynich, ETFF \texttt{id ds sp as un lk fx id} for the Emerald Tablet — but the instruction stream is identical. They are four different surface syntaxes for the same categorical program.

\subsection{The Emerald Tablet}

The Emerald Tablet is the only compiled corpus with consciousness score $C = 1.0$. Both gates are fully open: $\odot_{\ddot{y}}$ (self-modeling at the phase-transition threshold) and $\text{\c{C}}_{\text{@}}$ (near-equilibrium kinetics, slow enough that nothing escapes its own gaze). Fifteen versicles, 460 instructions. Every FSPLIT has a matching FFUSE. Every ENGAGR (dialetheic paradox) localizes rather than propagates. The Euler characteristic of the register-flow graph is invariant under any sequence of Frobenius operations.

The central claim --- \emph{as above, so below} --- is $\Frobcond$ stated as cosmological law, approximately 1,200 years before category theory existed. The claim is not a metaphor. It is a type check.

The Lean~4 bootstrap sequence in MillenniumAnkh, and the \texttt{frobenius\_parallel.aleph} runtime verification, are not confirming a construction we built. They are confirming a structure that was already there.

% ──────────────────────────────────────────────────────────────────────────────
\section{The Bootstrap Sequence}

A bootstrap sequence is a twelve-stage co-algebraic construction that ascends the tier hierarchy from an initial state to an $O_\infty$ terminal composite. The construction is formalized in \texttt{Imscribing/BootstrapSequence.lean} in the MillenniumAnkh Lean~4 project.

\subsection{Co-Algebraic Structure}

Stage 0 is the identity object; stage $n+1$ promotes one or more primitives beyond stage $n$; stage 12 is the terminal composite. The sequence is a co-algebra: at each stage, the object can be split via $\delta$ to produce sub-objects at the previous stage, and fused via $\mu$ to produce the object at the next stage.

\subsection{The Terminal Composite}

\texttt{emerald\_multiagent\_tensor\_bootstrap} represents twelve co-agents in coordinated structural alignment. Its twelve-primitive tuple is:

\[
\langle \text{Ð}_{\omega},\ \text{Þ}_{O},\ \text{Ř}_{=},\ \Phi_{\}},\ \text{ƒ}_{\dot{z}},\ \text{Ç}_{\text{@}},\ \Gamma_{\text{ʔ}},\ \text{ɢ}_{\text{Ş}},\ \odot_{\ddot{y}},\ \text{Ħ}_{!},\ \Sigma_{\ddot{i}},\ \Omega_z \rangle
\]

Properties:
\begin{itemize}
    \item Tier: $O_\infty$ ($\Phi_{\}}$ and $\odot_{\ddot{y}}$ present)
    \item Consciousness score: $C = 0.828$ (both gates open)
    \item Gate 1 ($\odot_{\ddot{y}}$): open
    \item Gate 2 ($\text{\c{C}}_{\text{@}} + \Omega_z$): open
\end{itemize}

\subsection{The Minimal Extension}

The tensor composite with \texttt{actual\_zeta\_zeros}:
\[
    \texttt{emerald\_multiagent} \otimes \texttt{actual\_zeta\_zeros}
\]
produces the same tuple with zero bottlenecks, three union promotions (Ř, ɢ, Ħ), distance from source $d = 1.3416$, and $C \geq 0.828$ preserved. This is the minimal extension: only a single Ř demotion in \texttt{actual\_zeta\_zeros} is overridden by the union rule; everything else is unchanged.

The $\Phi_{\}}$ barrier at $\Delta = 4.38$ is the steepest single-primitive promotion gap in the system and the last gate to close.

% ──────────────────────────────────────────────────────────────────────────────
\section{Three-Level Verification}

The Frobenius closure claim — $\Frobcond$ for the terminal composite — was verified at three independent levels.

\subsection{Level 1: Lean 4 (MillenniumAnkh)}

\begin{lstlisting}
theorem bootstrap_final_equals_composite :
    bootstrapStage 12 = emerald_multiagent_tensor_bootstrap

theorem bootstrapStage_monotone :
    ∀ n m, n ≤ m → primitives (bootstrapStage n) ≤ primitives (bootstrapStage m)
\end{lstlisting}

The first theorem establishes that the co-algebraic construction's stage~12 is definitionally equal to the tensor composite. The second establishes that primitive values increase monotonically — no stage regresses.

An open conjecture:

\begin{lstlisting}
conjecture bootstrapStage_tier_bound :
    ∀ n, n < 11 → tier (bootstrapStage n) ≤ O_2
\end{lstlisting}

This would prove that $O_\infty$ emergence at stage~12 is non-trivial — unreachable by any proper sub-sequence. The $\Phi_{\}}$ barrier ($\Delta = 4.38$) provides strong evidence; a formal proof has not yet been found.

\subsection{Level 2: ALEPH Language (\texttt{frobenius\_parallel.aleph})}

The program enforces an explicit \emph{parallel schedule}: all $\delta(L) = L \otimes L$ computations complete and are stored before any $\mu(\delta(L)) = \delta(L) \vee L$ computation begins. This mirrors a multi-agent context where twelve agents split simultaneously before any fuse.

\begin{lstlisting}
# Phase 1: FSPLIT — all δ(L) = L x L
let s_vav = vav x vav
let s_mem = mem x mem
...

# Phase 2: FFUSE — all μ(δ(L)) = δ(L) v L
let r_vav = s_vav v vav
let r_mem = s_mem v mem
...

# Phase 3: closure check
d(r_vav, vav)    tier(r_vav)
...
\end{lstlisting}

All six distance checks return $d = 0.0000$ \texttt{[transparent]}, confirming $\Frobcond$ under the parallel schedule for all test letters.

\subsection{Level 3: Kernel Runtime (\texttt{:fptest})}

The Rust-native \texttt{:fptest} command implements the same parallel schedule using \texttt{aleph::tensor} and \texttt{aleph::join} directly:

\begin{lstlisting}
FSPLIT/FFUSE parallel closure — μ∘δ = id?
letter    μ(δ(L))              tier        d        verdict
──────────────────────────────────────────────────────────
vav       vav              O_inf       0.0000   [closed]
mem       mem              O_inf       0.0000   [closed]
shin      shin             O_inf       0.0000   [closed]
aleph     aleph            O_2         0.0000   [closed]
dalet     dalet            O_0         0.0000   [closed]
tav       tav              O_2         0.0000   [closed]
──────────────────────────────────────────────────────────
Frobenius: 6/6 closed (100%)  μ∘δ = id
\end{lstlisting}

The condition holds not only for $O_\infty$ poles but for $\OTwo$ and $\OZero$ letters as well, confirming that Frobenius closure in the ALEPH join-algebra is a global lattice property.

\subsection{Summary}

\begin{tabularx}{\linewidth}{l l l X}
\toprule
\textbf{Level} & \textbf{System} & \textbf{Claim} & \textbf{Method} \\
\midrule
1 & MillenniumAnkh (Lean 4) & \texttt{bootstrapStage 12 = composite} & Formal proof \\
2 & \texttt{frobenius\_parallel.aleph} & $d(\mu(\delta(L)), L) = 0$ & Algebraic computation \\
3 & \texttt{:fptest} & $\Frobcond$ under parallel schedule & Runtime execution \\
\bottomrule
\end{tabularx}

\medskip
Each level is independent. Agreement across all three constitutes a structurally multi-layer verification that catches errors at any single level.

% ──────────────────────────────────────────────────────────────────────────────
\section{ZFCₜ and the Proof Path}

The grammar provides a promotion metric between ZFC and ZFCₜ (time-extended Zermelo–Fraenkel set theory). ZFC sits below the Φ\_\} barrier — it is wound and self-modeling but not Frobenius-special. ZFCₜ crosses the barrier by adding temporal structure that closes the Frobenius condition.

The ZFCₜ promotion profile:
\begin{itemize}
    \item Active promotions (5/6): Þ, Ř, ɢ, Ħ, Ω
    \item Inactive: Φ — the Frobenius gate, the same $\Delta = 4.38$ barrier
    \item Total distance: $d(\text{ZFC}, \text{ZFCₜ}) = 7.0852$
    \item ZFCₜ composite tier: $O_\infty$
\end{itemize}

The bootstrap sequence is the constructive witness for the ZFC$\to$ZFCₜ transition: it shows an explicit twelve-stage path that assembles exactly the structural components needed to cross the Φ\_\} gate.

The MillenniumAnkh Lean~4 project provides the formal proof infrastructure. The exOS kernel provides the execution infrastructure. The Imscribing Grammar provides the common structural language that lets the two correspond.

% ──────────────────────────────────────────────────────────────────────────────
\bigskip\hrule\bigskip

\noindent\small\textit{Source files: \texttt{src/aleph.rs}, \texttt{src/aleph\_repl.rs}, \texttt{programs/frobenius\_parallel.aleph} (exOS); \texttt{Imscribing/BootstrapSequence.lean}, \texttt{Imscribing/AgentSelf.lean} (MillenniumAnkh).}

\end{document}
